\documentclass[10pt,a4paper]{amsart}
\usepackage[margin=19mm]{geometry}
\usepackage{amsmath,amssymb,mathtools}
\usepackage[scr=rsfs,cal=euler]{mathalfa}
\usepackage{xfrac}
\usepackage[unicode,hidelinks,bookmarks=false]{hyperref}
\newcommand{\ie}{i.e.\ }
\newcommand{\cf}{cf.\ }
\newcommand{\resp}{respectively\ }
\newcommand{\Subsection}{\subsection}
\providecommand{\nobreakdash}{\nobreak\hbox{-}}
\setlength{\emergencystretch}{2em}
\multlinegap0pt
\usepackage{xcolor}
\usepackage{amssymb}
\usepackage{bm}
\usepackage[normalem]{ulem}
\usepackage{mathtools}
\usepackage{tikz}
\usepackage{float}
\newtheorem{lem}{Lemma}
\def\dif{{\rm d}}
\def\log{{\rm log}}
\title{On the volume of K-semistable Fano manifolds}
\author{Chi Li, Minghao Miao}
\date{Source: arXiv:2506.17420v3 (CC BY 4.0)}
\begin{document}
\maketitle

\noindent\emph{Source and licence.} On the volume of K-semistable Fano manifolds. Version arXiv:2506.17420v3 (CC BY 4.0), distributed by arXiv under the Creative Commons Attribution 4.0 International licence, http://creativecommons.org/licenses/by/4.0/, as recorded in the arXiv licence metadata for this exact version and shown on the abstract page https://arxiv.org/abs/2506.17420v3. Full licence notice: \texttt{licenses/arxiv-2506.17420-CC-BY-4.0.txt}.\par\smallskip

\noindent\textit{Editorial scope.} Counted unit: the article's lem lem (source0.tex lines 2619--2622). The article's complete original proof of it is delivered with it (source0.tex lines 2623--2688, one \texttt{proof} environment of 66 lines). No mathematics is written, repaired or re-derived: the statement and the proof are the article's own lines, in the article's own notation, and no retained line is substituted or re-flowed.  No further source-local context is needed: every cross-reference of every retained line resolves inside the two delivered spans.  Nothing was omitted from any retained span, so the package carries no omission record.  No retained line contains a citation, so the wrapper carries no bibliography and no bibliography entry.  No retained line contains a figure environment, an image-inclusion command or drawing code, so no image is delivered and the package declares no asset dependency.  Editorial formatting only, in addition: an AMS \texttt{amsart} preamble that restates the article's own declarations and macros used by the retained lines, namely the environments \texttt{lem} on separate counters, together with the standard packages those lines need.  The retained environments print in this document as Lemma 1 in order of appearance, whereas the article prints the same items with the numbering of its own full document.  The complete native source of the article as distributed in the arXiv e-print for this exact version is bundled unchanged as \texttt{sources/180231/source0.tex}.
\begin{lem}
For $3\le d\le n-2$, 
    $(2n^n-\phi(\tau_0))+\Psi(\tau_0)>0$. 
\end{lem}

\begin{proof}
    This is equivalent to:
    \begin{eqnarray*}
      \phi(\tau_0)-\Psi(\tau_0)=  \Phi(\tau_0)-\phi(\tau_0)< 2n^n. 
    \end{eqnarray*}
  
$r=n-d+1$, $x=2(n-d+1)+d$, $x-d=2r$, $x-n=2(n-d+1)+d-n=r+1$. 
\begin{eqnarray*}
&&\int_0^d z^{d-3}(d-z)(x_0-z)^{n-d+1}(n-z) \dif z\\
&=&\int_0^d (z-x_0+x_0)^{d-3}(d-x_0+x_0-z)(x_0-z)^{n-d+1}(n-x_0+x_0-z)\,  \dif z\\
&=&\int_0^d \sum_{j=0}^{d-3}\binom{d-3}{j}(-1)^j x_0^{d-3-j}(x_0-z)^{j+r} ((x_0-z)^2-(3r+1)(x_0-z)+2r(r+1))\, \dif z\\
&=&\sum_{j=0}^{d-3}\binom{d-3}{j}(-1)^j x_0^{d-3-j}\left[\frac{(x_0-z)^{j+r+3}}{j+r+3}-\frac{3r+1}{j+r+2}(x_0-z)^{j+r+2}+\frac{2r(r+1)}{j+r+1}(x_0-z)^{j+r+1}\right]_0^d. 
\end{eqnarray*}


\begin{eqnarray*}
    S&=&\frac{e^{1-\frac{d}{2}}}{2^{d-3}}\frac{1}{(d-3)!}\left[(4-d)\gamma(d-2, x)+2 x^{d-2}e^{-x}\right]\\
    &=&e^{1-\frac{d}{2}}2^{d-3}\left[(4-d)P(N\ge d-2)+2(d-2)P(N=d-2)\right]\\
    &\le& e^{1-\frac{d}{2}}2^{d-3}\left((4-d)\frac{2(d-1)}{d-2}+2(d-2)\right)P(N=d-2)\\
    &=&e^{1-\frac{d}{2}}2^{d-2}\frac{d}{d-2}P(N=d-2)\\
    &\le& \frac{d-1}{d-2}\frac{e}{\sqrt{2\pi d}}.
\end{eqnarray*}
\begin{eqnarray*}
    P(N=d-2)&=&\frac{x^{-(d-2)}}{(d-2)!}e^{-x}=\frac{x^{-d}}{d!}e^{-x}(d(d-1))\\
    &\le& x^{-(d-2)}e^{-x}\frac{d(d-1)}{d^d e^{-d}\sqrt{2\pi d}}\\
    &=& 2^{2-d} e^{d/2} \frac{d-1}{d}\frac{1}{\sqrt{2\pi d}}
\end{eqnarray*}
$\frac{\frac{x^{j+1}}{(j+1)!}}{\frac{x^{j}}{(j)!}}=\frac{x}{j+1}\le \frac{x}{d-1}=\frac{d}{2(d-1)}=r$
\begin{eqnarray*}
   P(N\ge d-2)\le \frac{P(N=d-2)}{1-r}=\frac{2(d-1)}{d-2}P(N=d-2). 
\end{eqnarray*}
\begin{eqnarray*}
  \mu_d(\tau) &=& \frac{\Phi(\tau_0)-\phi(\tau_0)}{2n^n}\\
   &=&\frac{2^{-(n-d+1)}}{2n^n}\sum_{j=0}^{n-d+1}\frac{n-d-j+2}{n-d+2}\frac{n!}{j!(n+1-j)!}d^{n-j}(2(n-d+1))^j\cdot\\
   &&\cdot [n^2-d^2+n+2d-nj]\\
   &=&\frac{2^{-(n-d+1)}n!}{2 n^n (d-3)!(n-d+2)!}\int_0^d z^{d-3}(d-z)(x-z)^{n-d+1}(n-z)dz.
\end{eqnarray*}

$\Delta=d(1-\frac{d-2}{n})=\frac{d(r+1)}{n}$, $a_k=\binom{n+1}{r-k}(1-q)^{n+1-r+k}q^{r-k}$
\begin{eqnarray*}
    &&\frac{\mu^r}{2n^n}\\
    &\le &\frac{(n+1)^n}{2n^n}\frac{1}{r+1}(1+\frac{r-1}{d})\sum_{j=0}^r (k+1)(k+r+1)\frac{(n+1)!}{(r-k)!(n+1-(r-k))!}(\frac{d}{n+1})^{d+k}(\frac{r}{n+1})^{r-k}2^{-k}\\
    &\le& \frac{e}{2}\frac{1}{r+1}(1+\epsilon)\sum_{j=0}^r(k+1)(k+r+1)\frac{1}{(r-k)!}r^{r-k}2^{-k}(1+\frac{r}{d})^{-d}\\
    &\le& \frac{e}{2(r+1)}e^{-r}\sum_{j=0}^r \frac{(k+1)(k+r+1)r^{r-k}}{2^k(r-k)!}e^{r\epsilon/2}(1+\epsilon)\\
    &\le& R_\infty e^{r(r+2)/(2d)}.
\end{eqnarray*}
\begin{eqnarray*}
    (1+\epsilon)^{-d}=\exp(-d \log(1+\epsilon))\le \exp(-d \epsilon+d\epsilon^2/2)=e^{-d\epsilon}e^{r^2/d}
\end{eqnarray*}

\begin{eqnarray*}
    \mu_3(\tau_0)&=&\frac{1}{2n^n}2^{-n+2}n\int_0^3(3-z)(x-z)^{n-2}(n-x+x-z)dz\\
    &=&\frac{2(n-5)(n-2)^n+(n-1/2)^{n-1}(n^2+5n-5)}{n^n(n+1)}\rightarrow 2 e^{-2}+e^{-1/2}=0.877201<1.
\end{eqnarray*}
\begin{eqnarray*}
    \frac{\mu^n}{\mu^{n+1}}=2\frac{(n+1)^n}{n^n}(n-d+3)
    \frac{\int_0^dz^{d-3}(d-z)((2n-d+2)-z)^{n-d+1}(n-z)dz}{\int_0^d z^{d-3}(d-z)(2n-d+4-z)^{n-d+2}(n+1-z)dz}.
\end{eqnarray*}
\begin{eqnarray*}
    \left(\frac{2n-d+4-z}{2n-d+2-z}\right)^{n-d+1}=\left(1+\frac{1}{n-\frac{d}{2}+1-\frac{z}{2}}\right)^{(n-\frac{d}{2}-\frac{z}{2}+1)+(\frac{z-d}{2})}. 
\end{eqnarray*}
\begin{eqnarray*}
    \mu_{n-2}(\tau_0)&=&\frac{(n-2)^{n-2}(40-66n-126 n^2+115 n^3)}{16n^n(1+n)}\rightarrow e^{-2}\frac{115}{16}\sim 0.9728<1. 
\end{eqnarray*}

\end{proof}

\end{document}
